\subsubsection{加法作为投影的碰撞谱：10 状态同步核 vs.\ 40 状态真实输入核}\label{app:add-collision-spectrum}

\paragraph{加法 = 升维到 \texorpdfstring{$\{0,1\}^2$}{\{0,1\}^2} 再投影回 \texorpdfstring{$X_\infty$}{X\_infty}}
在稳定地址层，Zeckendorf 加法可写为
\[
c\oplus d=\Fold_\infty(c+d),
\]
其中 \((c+d)_j:=c_j+d_j\in\{0,1,2\}\) 为逐位和；\(\Fold_\infty\) 为把 \(\{0,1,2\}\) 归一化回 \(\{0,1\}\) 的在线归一化门，其可由 10 状态同步核 \(K_{\mathrm{sync}}\) 实现（表 \ref{tab:sync-kernel-A-compare}）。
因此“加法作为一次投影”的机制可被视为：
\[
\{0,1\}^{\ZZ}\times\{0,1\}^{\ZZ}\xrightarrow{\ (x,y)\mapsto x+y\ }\{0,1,2\}^{\ZZ}\xrightarrow{\ K_{\mathrm{sync}}\ }\{0,1\}^{\ZZ}\supseteq X_\infty,
\]
即先把输入升维到两条比特流（本体场），再由归一化核投影回稳定地址外表。

\paragraph{碰撞矩与空间损失谱常数}
把上式视为一个确定性转导器（输入字母表 \(\{0,1\}^2\)，输出字母表 \(\{0,1\}\)）。
对窗口长度 \(m\) 与输出块 \(y\in\{0,1\}^m\)，令
\[
d^{\mathrm{add}}_m(y):=\#\{(x,x')\in(\{0,1\}^m)^2:\ \Fold_m(x+x')=y\},
\qquad
S^{\mathrm{add}}_q(m):=\sum_{y\in\{0,1\}^m}d^{\mathrm{add}}_m(y)^q\quad(q\ge 2).
\]
则 \(S^{\mathrm{add}}_q(m)\) 的指数增长率由一族有限非负整数 moment-kernel \(A^{\mathrm{add}}_q\) 控制，记
\[
r^{\mathrm{add}}_q:=\rho\!\left(A^{\mathrm{add}}_q\right),
\qquad
S^{\mathrm{add}}_q(m)=\Theta\!\left((r^{\mathrm{add}}_q)^m\right).
\]
尤其当 \(q=2\) 时，有概率解释：若从 \(\{0,1\}^{2m}\) 上均匀独立抽取两对输入 \((x_1,x'_1),(x_2,x'_2)\)，则
\[
\mathbb{P}\bigl[\Fold_m(x_1+x'_1)=\Fold_m(x_2+x'_2)\bigr]=\frac{S^{\mathrm{add}}_2(m)}{16^m},
\qquad
h^{\mathrm{add}}_2=\log\!\left(\frac{16}{r^{\mathrm{add}}_2}\right).
\]

\paragraph{真实输入约束的 40 状态提升}
若要求两条输入流本身为合法 Zeckendorf 轨道（各自禁止相邻 \(11\)），则可把该约束并入状态，得到 40 状态真实输入同步核 \(K_{\mathrm{real}}\)（表 \ref{tab:sync-kernel-A-compare}）。
在该核上对同一“加法投影”做相同的碰撞 moment 编译，得到真实输入下的 \(r^{\mathrm{add,real}}_q\)。

\begin{remark}[常数级差异（可复现）]\label{rem:add-collision-spectrum-10-vs-40}
用仓库脚本 \texttt{scripts/exp\_add\_collision\_spectrum\_10\_vs\_40.py}（基于对称商/系数提取的 histogram-DP moment-kernel 实现）可复现计算：
\(\;r^{\mathrm{add}}_2\approx 8.743\;\)、
\(\;r^{\mathrm{add}}_3\approx 19.851\;\)，
而在真实输入约束下
\(\;r^{\mathrm{add,real}}_2\approx 3.999\;\)。
对照表见表 \ref{tab:add-collision-spectrum-10-vs-40}。
\input{sections/generated/tab_add_collision_spectrum_10_vs_40}

该结论给出一个直接、可量化的“40 状态核的独立价值”锚点：把输入合法性约束提升为核几何后，“加法作为投影”的空间侧碰撞常数发生显著改变（常数因子级下降），从而其空间损失谱不再可由 10 状态归一化器的无约束口径替代。
\end{remark}

